\documentclass[11pt]{article}
\usepackage[margin=2.5cm]{geometry}
\usepackage{hyperref}
\usepackage{xcolor}
\usepackage{enumitem}
\usepackage{booktabs}
\usepackage{parskip}

\title{\textbf{Response to Reviewers}\\[0.5em]
  \large GraphyloVar: Predicting the impact of non-coding variants\\
  using a multi-species sequence model\\[0.5em]
  \normalsize Bioinformatics, Manuscript BIOINF-2025-2871}
\author{Dongjoon Lim and Mathieu Blanchette}
\date{}

\begin{document}
\maketitle

We thank the reviewers for their reading and constructive feedback.
We addressed all comments and the revised manuscript marks all new or
changed text in red.

\bigskip
\hrule
\section*{Reviewer 1}

\subsection*{R1.1 GCN interpretability and species-level analysis}

\begin{quote}
\itshape
The manuscript proposes the use of GCN to integrate multi-species
sequence information in order to capture cross-species signals.
However, the current presentation lacks a detailed analysis of this
component.  The author is encouraged to include additional
investigations into the correlations among multi-species sequences
and to enhance the interpretability of the sequence representations
learned by the GCN module.
\end{quote}

We added a new section (Section~3.6, GCN Interpretability and Species
Importance) that provides a detailed additive perturbation analysis of
the GCN component.  For each of the 58 extant placental mammal species,
we replaced all other alignment rows with a gap token and measured the
decrease in nucleotide cross-entropy relative to an all-gap baseline,
evaluated on 500,000 held-out variants from chromosomes~13 to~22.  This
quantifies the isolated information content of each species.

The results show that 53 of 58 extant species contribute positive signal.
Human provides the largest contribution ($\Delta$CE $= 7.1$).  Four
diverged primates (Orangutan, Green monkey, Gorilla, Gibbon;
$\Delta$CE $= 0.62$ to $1.50$) form a secondary cluster.  Chimpanzee
ranks 11th ($\Delta$CE $= 0.062$) despite being the closest relative of
humans because its sequence is nearly identical to the Human row, leaving
no additional marginal signal (leave-one-out marginal $\Delta$CE
$\approx -2.4 \times 10^{-6}$).  A conditional perturbation analysis in
which the Human row is also masked confirms this: Chimpanzee becomes the
top non-human contributor in that setting.  The SE-attention gate values
learned by the model are consistent with these perturbation rankings,
showing the model learns to weight species according to their
phylogenetic informativeness rather than treating them equally.

Section~3.6 added; Figure~6 (species importance bar chart, clade-colored,
log scale) and Supplementary Figures~S1 (phylogenetic tree ablation) and
S2 (model component ablation) added.

\bigskip

\subsection*{R1.2 Justification for 65 bp context window}

\begin{quote}
\itshape
In the Methods section, a fixed sequence fragment length of 65 bp is
adopted for feature extraction, yet no sufficient justification for
this parameter choice is provided.  To strengthen the methodological
rigor, the author should consider conducting systematic experiments to
demonstrate the optimality or robustness of this setting.
\end{quote}

We trained models with three context-window settings: flank $\in \{16,
32, 100\}$ bp (input sequences of 33, 65, and 201 bp).  All models used
the same training procedure on chromosomes~1 to~10 with chromosomes~11
to~12 as validation.  Flank~$= 32$ (65~bp) achieves the highest overall
held-out AUROC (0.625) among the three settings, and also achieves the
highest AUROC on all 13 MPRA benchmark datasets after fine-tuning,
supporting this choice for downstream regulatory prediction tasks.
These results are reported in Supplementary Table~S1.

Section~3.5 and Supplementary Table~S1 added.

\bigskip

\subsection*{R1.3 Quantitative evidence for complementarity}

\begin{quote}
\itshape
In the section ``Complementarity of GraphyloVar Scores with Existing
Methods,'' the author claim that GraphyloVar identifies unique functional
variant information that existing approaches fail to capture.  However,
this assertion is not supported by corresponding quantitative or
qualitative experimental evidence.  Additional analyses are needed to
substantiate the claimed complementarity and to demonstrate the unique
contributions of GraphyloVar.
\end{quote}

We added quantitative evidence in Section~3.4.  We computed a
parameter-free z-score ensemble of GraphyloVar and CADD on the full
$\sim$149M held-out variants.  Both scores were independently z-normalized;
because CADD assigns higher scores to more deleterious variants (i.e.,
rarer ones), its z-score was negated to align direction before averaging.
The ensemble achieves AUROC~$= 0.6442$, an improvement of $+0.020$ over
GraphyloVar alone ($p < 10^{-15}$, DeLong test on 500,000 held-out
variants).  Bootstrap 95\% confidence intervals ($B = 1{,}000$) for the
ensemble and GraphyloVar alone do not overlap.  The low Spearman correlation
between GraphyloVar and CADD shown in Figure~4, together with this AUROC
gain, confirms that GraphyloVar captures information that existing methods
do not.

Section~3.4 expanded; Figure~5 added.

\bigskip

\subsection*{R1.4 Comparison with Evo2 and GPN-Star}

\begin{quote}
\itshape
Several existing approaches leverage multi-species sequence information
to predict the functional effects of non-coding variants, including Evo2
and GPN-Star.  These models have demonstrated strong performance in
integrating evolutionary conservation with functional annotations.  To
more comprehensively assess the performance, novelty, and practical
advantages of GraphyloVar, it is recommended that the author include a
quantitative comparison with Evo2 and GPN-Star in the experimental
evaluation section.\\
{[1]}~Genome modeling and design across all domains of life with Evo 2.\\
{[2]}~Predicting functional constraints across evolutionary timescales
with phylogeny-informed genomic language models.
\end{quote}

GPN-Star was added as a baseline in the common vs.\ rare variant AUROC
benchmark (Figure~2A).  Due to the computational cost of scoring
$\sim$149M variants with GPN-Star, Spearman correlation with observed MAF
was not computed for GPN-Star, so it does not appear in Table~1.  GPN-Star
was also not included in the MPRA evaluation because fine-tuning GPN-Star
on MPRA data was outside the scope of this work.

Evo2 requires GPU compute capability 8.9 or higher (Ada or Hopper
architecture).  Our hardware consists of NVIDIA Quadro RTX 6000 GPUs
(compute capability 7.5), which is below this requirement.  We state
this explicitly in Methods Section~2.4 and note Evo2 as an important
future comparison in the Discussion.

Methods Section~2.4 and Discussion updated with GPN-Star paragraph and
Evo2 hardware justification.

\bigskip

\subsection*{R1.5 Code and data availability}

\begin{quote}
\itshape
The author provides a code repository link in the manuscript; however,
the repository does not include the complete implementation, nor does
it contain the datasets used in the experiments.  To ensure the
reproducibility of the study, it is strongly recommended that the
author release the full source code along with the corresponding
datasets, and clearly specify the procedures for accessing them as
well as their usage licenses.
\end{quote}

A Data and Code Availability section was added at the end of the
manuscript.  It provides the GitHub and McGill PUB URLs for source code, trained model
weights (flank~$= 32$ checkpoint), and preprocessing scripts, released
under the MIT License.  It also provides the dbGaP accession number
(phs000964) for the TOPMed whole-genome sequencing data, the UCSC URL for
the 100-way vertebrate alignments, and pointers to the original
publications for each MPRA dataset.

Data and Code Availability section added.

\bigskip

\subsection*{R1.6 Dataset descriptions and figure label consistency}

\begin{quote}
\itshape
On Page 6, Line 10, the manuscript mentions the use of ``saturation
mutagenesis of disease-associated regulatory elements [Kircher et al.,
2019],'' yet no direct comparison or explicit presentation of results
based on this dataset is provided in the main text.  In addition,
Figure 3 and Figure 5 includes an annotation referring to the
``GWAScatalog'' for 3' UTRs of genomic regulatory contexts, but the
database is not described elsewhere in the manuscript.  The author
should supply the relevant dataset descriptions and ensure consistency
between cited resources and the results presented.
\end{quote}

All 13 MPRA datasets are now explicitly described in Methods Section~2.5,
grouped by the three categories shown in Figure~3.  Kircher et al.\ (2019)
provided functional scores for all possible single-nucleotide variants
(all three alternative alleles at each position) within disease-associated
regulatory elements, forming part of the genome-wide
analyses group along with Abell et al.\ (2022) and McAfee et al.\ (2023).
For the 3'~UTR group, Griesemer et al.\ (2021) used annotations from the
GWAS Catalog to identify causal 3'~UTR variants, and Schuster et al.\
(2023) characterized 3'~UTR variants associated with prostate cancer.
The label used in Figures~3 and~5 refers to the GWAS Catalog as the
annotation source for that dataset, and its meaning is now explained
in the Methods text.

Methods Section~2.5 rewritten to describe all 13 MPRA datasets.

\bigskip

\subsection*{R1.7 MAF correlation reported only in text}

\begin{quote}
\itshape
The section ``Correlation with Observed Minor Allele Frequencies''
reports result only in textual form.  To improve the clarity and
credibility of the conclusions, it would be beneficial to include
relevant supplementary figures or tables.
\end{quote}

The Spearman correlation values are reported in Table~1 (Section~3.2).
The table includes GraphyloVar and five baselines (GPN-MSA, Enformer,
PhyloP, CADD, PhastCons).  GPN-Star Spearman correlation was not computed
as described in R1.4 above.

Table~1 added to Section~3.2.

\bigskip

\subsection*{R1.8 Figure~1 architecture inconsistency and white text}

\begin{quote}
\itshape
In Figure 1 (Page 4), the features after Mean Pooling are shown as an
input to the ``Scaling species with Attention'' module.  Since mean
pooling typically produces a one-dimensional vector, this appears
inconsistent with the subsequent GCN module, which seems to require
high-dimensional feature inputs.  The author should clarify the
rationale behind this architectural design or revise the figure to
accurately reflect the implemented model.  Additionally, there is
unexplained white text near the ``Concatenate Left + Right'' in the
figure, the author should verify its validity and correct or remove it
as appropriate.
\end{quote}

Methods Section~2.2 and the Figure~1 caption were updated to clarify the
pooling operation.  After processing both the forward and reverse complement
strands independently through two Transformer layers each, the output for
each strand is mean-pooled over the 65 input positions, producing a
32-dimensional representation per strand.  The forward and reverse
complement strand representations are concatenated into a 64-dimensional
feature vector per species (32 from the forward strand and 32 from the
reverse complement strand), which is then passed to the SE block and GCN.
This is consistent with the ``Concatenate Left + Right'' operation shown
in Figure~1.

The white text near ``Concatenate Left + Right'' in the original figure is
a rendering artifact from the figure generation tool and does not correspond
to any model component.

Methods Section~2.2 and Figure~1 caption updated.

\bigskip

\subsection*{R1.9 Figure~2 caption label error}

\begin{quote}
\itshape
In Figure 2 (Page 7), the figure caption refers to ``GraphyloVar Score
(dark purple) and GraphyloVar Score (purple)'' which appears redundant
or incorrect.  The author should carefully check the caption text,
legend colors, and corresponding labels for accuracy and consistency.
\end{quote}

The Figure~2 caption was rewritten to remove color references.  The
updated caption describes each panel directly without redundant or
confusing color labels, since the figure legend already identifies each
method.

Figure~2 caption rewritten.

\bigskip

\subsection*{R1.10 Naming consistency: GPN-MSA}

\begin{quote}
\itshape
There is inconsistency in the naming of the proposed method throughout
the manuscript (e.g., ``GPN-MSA'' vs.\ ``GPNMSA'').  The author should
standardize the nomenclature across the entire text and explicitly
clarify whether the hyphen should be retained.
\end{quote}

We use ``GPN-MSA'' (with hyphen) consistently throughout the manuscript.
We verified this with \texttt{grep}; no instances of ``GPNMSA'' without
a hyphen remain in any of the source files.

Global find-replace applied.

\bigskip
\hrule
\section*{Reviewer 2}

\subsection*{R2.1 Held-out data definition and potential leakage}

\begin{quote}
\itshape
Held-out data definition: The manuscript states that variants are
``held out,'' but it is unclear whether this refers only to
variant-level exclusion or also to separation at the level of genomic
regions.  This should be stated explicitly to rule out potential
leakage.  If possible, an evaluation of region-based hold out would
strengthen the work.
\end{quote}

The holdout strategy is explicitly stated in Section~3.1.  Chromosomes~1
to~10 were used for training, chromosomes~11 to~12 for validation, and
chromosomes~13 to~22 exclusively for the held-out test evaluation.  This
chromosome-level partition ensures that no variant or genomic position
from the test chromosomes appears anywhere in training or validation.  The
test set contains approximately 149~million SNVs from the TOPMed
whole-genome sequencing cohort.  This chromosome-level split provides
separation at both the variant level and the genomic region level, because
entire chromosomes are withheld.

For MPRA fine-tuning evaluation, leave-one-out cross-validation was used
within each dataset group, as described in Section~3.3.

Section~3.1 first paragraph updated.

\bigskip

\subsection*{R2.2 Discussion of PHACT and PHACTboost}

\begin{quote}
\itshape
Prior work that incorporates phylogenetic information, such as PHACT
(Kuru et al.\ 2022) and PHACTboost (Kuru et al., 2024), should be
discussed to better link GraphyloVar within the phylogeny-aware variant
prediction literature, even if the modeling approach differs or the
work focuses on coding regions.
\end{quote}

A comparison paragraph with PHACT (Kuru et al., 2022) and PHACTboost
(Kuru et al., 2024) was added to the Discussion.  PHACT and PHACTboost
use protein-level features including amino acid conservation, protein
structure, and functional domain annotations, and are designed specifically
for missense variants in coding regions.  GraphyloVar targets both coding
and non-coding variants using only multi-species whole-genome alignments
and does not use any protein-level features.  The two approaches are
therefore complementary: PHACT provides high accuracy for missense variant
interpretation while GraphyloVar provides genome-wide coverage across all
variant types.  We note in the Discussion that combining them in an
ensemble for coding variants would be an interesting direction for future
work.

Discussion paragraph added.

\bigskip

\subsection*{R2.3 Comparison breadth relative to dbNSFP}

\begin{quote}
\itshape
The dbNSFP hosts 36 algorithms' prediction scores.  Comparing to a
larger set of tools should be possible, and if not, an explicit
justification for their exclusion should be included.
\end{quote}

A justification paragraph was added to Methods Section~2.4.  We selected
baselines representing methodologically distinct families: conservation
metrics (PhyloP, PhastCons), integrative scoring (CADD), sequence-based
deep learning (Enformer), and MSA-aware models (GPN-MSA, GPN-Star).
Running all 36 dbNSFP tools was outside our computational scope; many
share underlying conservation or annotation features, so the selected
families collectively cover the major methodological landscape of variant
effect prediction without redundancy.

Methods Section~2.4 updated with method selection justification.

\bigskip

\subsection*{R2.4 Loss function clarification}

\begin{quote}
\itshape
The allele-frequency head outputs a categorical distribution, yet
binary cross-entropy is reported as the loss.  This requires
clarification.
\end{quote}

We corrected this in Methods Section~2.3.  The allele frequency head
uses categorical cross-entropy loss because it predicts a probability
distribution over five allele states (A, C, G, T, gap).  The SNP
probability head uses binary cross-entropy loss because polymorphism
status is a binary label.  The original text incorrectly stated binary
cross-entropy for both heads.

Methods Section~2.3 corrected.

\bigskip
\hrule
\bigskip

We thank the reviewers again for their careful reading and comments.
We believe the revised manuscript fully addresses all the points raised.

\end{document}
